\documentclass[a4paper]{article}

% --- Основные пакеты ---
\usepackage{graphicx}     % картинки
\usepackage{geometry}     % поля страницы
\usepackage{hyperref}     % гиперссылки
\usepackage{amsmath}      % математика
\ifPDFTeX
    \usepackage{amssymb}  % математические символы
\else
    \usepackage{unicode-math} % масштабируемые OpenType-матшрифты
\fi
\usepackage{float}        % [H] позиционирование фигур
\usepackage{iftex}        % определение движка (pdfLaTeX / XeLaTeX / LuaLaTeX)
\ifPDFTeX
    \usepackage[T2A]{fontenc} % кириллица
    \usepackage[utf8]{inputenc} % кодировка UTF-8
    \usepackage{anyfontsize}  % точные размеры шрифтов без подстановки
\else
    \usepackage{fontspec}     % системные шрифты с поддержкой кириллицы
    \setmainfont{Times New Roman}
    \setsansfont{Arial}
    \setmonofont{Courier New} % Consolas нет в MS core fonts на Linux — \texttt{} выводился бы пустым
\fi
\usepackage[russian]{babel} % многоязычность
\usepackage{caption}      % настройка подписей
\usepackage{booktabs}     % линейки в таблицах
\usepackage[fontsize=13pt]{fontsize} % Размер шрифта

% --- Настройка страницы ---
\geometry{top=1.5cm, bottom=1.5cm, left=3cm, right=1.5cm}

% --- Картинки: относительный путь к code/results относительно report/ (кавычки обязательны из-за пробелов в пути) ---
\newcommand{\img}[2][]{\includegraphics[#1]{{"../code/results/#2"}.png}}

% --- Настройка гиперссылок ---
\hypersetup{
    colorlinks,
    linkcolor=black,
    urlcolor=blue,
    citecolor=black
}

% --- Настройка подписей ---
\renewcommand{\figurename}{Рис.}
\captionsetup{labelsep=period}

% Заголовки разделов и подразделов без нумерации
\newcommand{\mainsection}[1]{%
    \phantomsection
    \section*{\centering\fontsize{16pt}{0pt}\selectfont\textbf{#1}}%
}
\newcommand{\subsectiontext}[1]{%
    \section*{\fontsize{13pt}{0pt}\selectfont\textbf{#1}}%
}

\begin{document}

% --- Титульная страница ---
\centering
Министерство науки и высшего образования Российской Федерации\\
Федеральное государственное автономное образовательное учреждение\\
высшего образования\\
«Южно-Уральский государственный университет\\
(национальный исследовательский университет)»\\
Институт естественных и точных наук

\setlength{\parskip}{7.5cm}

Отчет\\
Лабораторная работа №2.\\
Методы численной оптимизации.

\setlength{\parskip}{5.9cm}

\raggedleft
\begin{minipage}{0.46\textwidth}
\raggedright
Выполнил: студент группы ЕТ-310\\
Тельнов С.М.\\
Проверил: преподаватель кафедры ВМ\\
Стариков Я.Е.
\end{minipage}

\setlength{\parskip}{5.84cm}

\centering
Челябинск — 2026

\setlength{\parskip}{0pt plus 1pt}
\newpage

% --- Содержание ---
\centering
\textbf{Краткое содержание}
\begin{itemize}
    \setlength{\parskip}{0.4cm}
    \item \hyperref[sec:theory]{\text{Теоретическая часть} \dotfill \pageref{sec:theory}}
    \item \hyperref[sec:practice]{\text{Практическая часть} \dotfill \pageref{sec:practice}}
    \item \hyperref[sec:conclusion]{\text{Вывод} \dotfill \pageref{sec:conclusion}}
\end{itemize}
\newpage

\setlength{\parindent}{1.5em}

\raggedright

% --- Теоретическая часть ---
\label{sec:theory}
\mainsection{Теоретическая часть}

\subsectiontext{1. Постановка задачи и структура программного комплекса}

Решается задача безусловной минимизации скалярной функции $f: \mathbb{R}^n \to \mathbb{R}$:
\[
    \vec{\theta}^{\,*} = \arg\min_{\vec{\theta}} f(\vec{\theta}),
\]
а также задача с ограничениями типа равенств и неравенств (блок 5). Все методы реализованы на языке Python (среда управления зависимостями \texttt{uv}, библиотеки \texttt{numpy} и \texttt{matplotlib}, линтер \texttt{ruff}) и разбиты на модули:

\begin{itemize}
    \item \texttt{objectives.py} — целевые функции и численное дифференцирование;
    \item \texttt{optimizers.py} — градиентные методы первого и второго порядка, одномерный поиск шага;
    \item \texttt{zeroth\_order.py} — методы нулевого порядка (Нелдер–Мид, рой частиц);
    \item \texttt{least\_squares.py} — нелинейный метод наименьших квадратов;
    \item \texttt{session\_task.py} — прикладная задача с ограничениями, мультизапуск и кластеризация;
    \item \texttt{main.py} — запуск пяти блоков расчётов и построение всех графиков.
\end{itemize}

Ключевые классы и их назначение:

\textbf{Класс \texttt{TargetFunction}} (файл \texttt{objectives.py}) — базовый класс скалярной целевой функции. Хранит ссылку на функцию $f(\vec{\theta})$, а также аналитические градиент и гессиан, заданные при создании. При отсутствии аналитических выражений методы \texttt{grad} и \texttt{hess} автоматически переходят на численные оценки: градиент вычисляется центральными конечными разностями
\[
    \frac{\partial f}{\partial \theta_i}(\vec{\theta}) \approx
    \frac{f(\vec{\theta} + h \vec{e}_i) - f(\vec{\theta} - h \vec{e}_i)}{2h},
    \qquad h = 10^{-5},
\]
а численный гессиан получается как якобиан численного градиента (вспомогательная функция \texttt{jacobian}, $h = 10^{-4}$). Класс также хранит начальную точку $\vec{\theta}^{\,0}$ и параметры области построения изолиний. От него наследуются:
\begin{itemize}
    \item \texttt{Rosenbrock(n)} — функция Розенброка $f(\vec{\theta}) = \sum_{i=1}^{n-1}\left[100\,(\theta_{i+1} - \theta_i^2)^2 + (1 - \theta_i)^2\right]$ с единственным минимумом $f^* = 0$ в точке $(1, \dots, 1)^T$;
    \item \texttt{Himmelblau} — функция Химмельблау $f(\vec{\theta}) = (\theta_1^2 + \theta_2 - 11)^2 + (\theta_1 + \theta_2^2 - 7)^2$ с четырьмя минимумами;
    \item \texttt{ExamSession} — целевая функция прикладной задачи распределения времени (блок 5).
\end{itemize}

\textbf{Класс \texttt{LineSearch}} (файл \texttt{optimizers.py}) — одномерная минимизация шага $\alpha$ вдоль луча $\vec{\theta} + \alpha \vec{d}$ методом золотого сечения. Интервал неопределённости сначала расширяется вдвое, пока значение $\varphi(\alpha) = f(\vec{\theta} + \alpha \vec{d})$ убывает, затем сжимается с коэффициентом $r = \frac{3 - \sqrt{5}}{2} \approx 0{,}618$ до достижения точности $10^{-8}$.

\textbf{Класс \texttt{Optimizer}} — базовый класс итерационного метода спуска. Его метод \texttt{run} реализует общий цикл: вычисление градиента, проверку критерия остановки $\|\nabla f\| < \texttt{tol}$, выполнение шага через абстрактный метод \texttt{\_move} и запись истории. После завершения доступны история приближений \texttt{history\_}, значения функции \texttt{f\_history\_} и нормы градиента \texttt{grad\_norms\_}. Конкретный алгоритм задаётся переопределением \texttt{\_move}, сброс внутреннего состояния — переопределением \texttt{\_reset}. От этого класса наследуются все градиентные методы (\texttt{SteepestDescent}, \texttt{CoordinateDescent}, \texttt{NesterovMomentum}, \texttt{Adam}, \texttt{Adagrad}, \texttt{RMSProp}, \texttt{Newton}, \texttt{BFGS}), что позволяет единообразно сравнивать их траектории.

\textbf{Классы нулевого порядка} (файл \texttt{zeroth\_order.py}): \texttt{NelderMead} — симплекс-метод с сохранением истории вершин симплекса; \texttt{PSO} — метод роя частиц с сохранением снимков позиций и скоростей частиц на каждой итерации.

\textbf{Класс \texttt{LeastSquaresProblem}} (файл \texttt{least\_squares.py}) — постановка задачи МНК: хранит векторную функцию невязок $\vec{r}(\vec{\theta})$, её якобиан $J$ (аналитический или конечные разности), целевую функцию $S(\vec{\theta}) = \vec{r}^{\,T}\vec{r}$ и градиент $\nabla S = J^T \vec{r}$.

\subsectiontext{2. Градиентные методы первого порядка}

Все методы блока используют только значения функции и её первых производных. Градиент вычисляется двумя способами — аналитически и численно (центральные разности), что позволяет проверить корректность реализации дифференцирования.

\textbf{Наискорейший градиентный спуск} (\texttt{SteepestDescent}). Итерация:
\[
    \vec{\theta}^{\,k+1} = \vec{\theta}^{\,k} - \alpha_k \nabla f(\vec{\theta}^{\,k}),
    \qquad
    \alpha_k = \arg\min_{\alpha \ge 0} f\bigl(\vec{\theta}^{\,k} - \alpha \nabla f(\vec{\theta}^{\,k})\bigr),
\]
где одномерный минимум находится классом \texttt{LineSearch} (золотое сечение). Если задан параметр \texttt{lr}, одномерный поиск отключается и выполняется спуск с фиксированным множителем $\text{lr}$:
\[
    \vec{\theta}^{\,k+1} = \vec{\theta}^{\,k} - \text{lr}\, \nabla f(\vec{\theta}^{\,k}).
\]
Такой вариант дешевле по числу вычислений $f$, но шаг должен удовлетворять условию устойчивости $\text{lr} < 2/L$ ($L$ — константа Липшица градиента), поэтому на плохо обусловленных функциях он либо расходится, либо сходится медленно. В расчётах $\text{lr} = 2\cdot10^{-3}$.

\textbf{Покомпонентный спуск} (\texttt{CoordinateDescent}). Координаты обновляются по одной в цикле Гаусса–Зейделя, причём градиент пересчитывается после каждого шага компоненты:
\[
    \theta_i \leftarrow \theta_i - \alpha_i \frac{\partial f}{\partial \theta_i}(\vec{\theta}),
    \qquad
    \alpha_i = \arg\min_{\alpha \ge 0} f\bigl(\vec{\theta} - \alpha\, \text{grad}_i f(\vec{\theta})\, \vec{e}_i\bigr).
\]
Реализованы два варианта обхода осей: последовательный ($\theta_1 \to \theta_2 \to \dots$) и со случайным перемешиванием порядка координат на каждой итерации (\texttt{shuffle=True}, генератор с фиксированным зерном). Как и у наискорейшего спуска, при заданном \texttt{lr} одномерный поиск заменяется фиксированным шагом $\theta_i \leftarrow \theta_i - \text{lr}\, \partial f / \partial \theta_i$ с тем же $\text{lr} = 2\cdot10^{-3}$.

\textbf{Метод Нестерова} (\texttt{NesterovMomentum}) — спуск с ускорением, градиент вычисляется в опережающей (look-ahead) точке:
\[
    \vec{v}^{\,k+1} = \mu\, \vec{v}^{\,k} - \text{lr}\, \nabla f\bigl(\vec{\theta}^{\,k} + \mu\, \vec{v}^{\,k}\bigr),
    \qquad
    \vec{\theta}^{\,k+1} = \vec{\theta}^{\,k} + \vec{v}^{\,k+1},
\]
параметры: $\mu = 0{,}95$, $\text{lr} = 10^{-3}$.

\textbf{Метод Adam} (\texttt{Adam}) — адаптивный шаг с экспоненциальными скользящими средними градиента и его квадрата:
\[
    \vec{m}_t = \beta_1 \vec{m}_{t-1} + (1-\beta_1)\, \nabla f_t,
    \qquad
    \vec{v}_t = \beta_2 \vec{v}_{t-1} + (1-\beta_2)\, (\nabla f_t)^2,
\]
с поправками на смещение $\hat{\vec{m}}_t = \vec{m}_t / (1-\beta_1^t)$, $\hat{\vec{v}}_t = \vec{v}_t / (1-\beta_2^t)$ и шагом
\[
    \vec{\theta}^{\,t+1} = \vec{\theta}^{\,t} - \text{lr}\, \frac{\hat{\vec{m}}_t}{\sqrt{\hat{\vec{v}}_t} + \varepsilon},
    \qquad
    \beta_1 = 0{,}9,\ \beta_2 = 0{,}999,\ \varepsilon = 10^{-8},\ \text{lr} = 0{,}05.
\]

\textbf{Метод Adagrad} (\texttt{Adagrad}) — адаптивный шаг с накапливаемой суммой квадратов компонент градиента:
\[
    \vec{G}_t = \vec{G}_{t-1} + (\nabla f_t)^2,
    \qquad
    \vec{\theta}^{\,t+1} = \vec{\theta}^{\,t} - \text{lr}\, \frac{\nabla f_t}{\sqrt{\vec{G}_t} + \varepsilon},
\]
все операции покомпонентные, $\text{lr} = 1$, $\varepsilon = 10^{-8}$. Редко обновляемые координаты получают больший шаг, но монотонный рост $\vec{G}_t$ постепенно глушит шаг по всем координатам.

\textbf{Метод RMSProp} (\texttt{RMSProp}) снимает проблему затухания шага, заменяя накапливаемую сумму экспоненциальным скользящим средним:
\[
    \vec{v}_t = \beta\, \vec{v}_{t-1} + (1-\beta)\, (\nabla f_t)^2,
    \qquad
    \vec{\theta}^{\,t+1} = \vec{\theta}^{\,t} - \text{lr}\, \frac{\nabla f_t}{\sqrt{\vec{v}_t} + \varepsilon},
\]
$\beta = 0{,}999$, $\text{lr} = 0{,}01$, $\varepsilon = 10^{-8}$. Окно усреднения ограничено, поэтому шаг не стремится к нулю; при постоянном $\text{lr}$ вблизи минимума (где $\vec{v}_t \approx (\nabla f_t)^2$) итерации превращаются в блуждание с амплитудой $\sim \text{lr}$.

\subsectiontext{3. Методы второго порядка}

\textbf{Классический метод Ньютона} (\texttt{Newton}):
\[
    \vec{\theta}^{\,k+1} = \vec{\theta}^{\,k} - H_f^{-1}(\vec{\theta}^{\,k})\, \nabla f(\vec{\theta}^{\,k}).
\]
Гессиан вычисляется аналитически либо численно. Система решается методом наименьших квадратов (\texttt{lstsq}); если направление не является направлением спуска ($\nabla f^T \vec{p} \ge 0$) или вырождено, выполняется откат на антиградиент, а если полный шаг не уменьшает функцию — одномерный поиск вдоль направления Ньютона.

\textbf{Квазиньютоновский метод BFGS} (\texttt{BFGS}) обходится без точного гессиана, итерационно аппроксимируя обратный гессиан $B_k \approx H_k^{-1}$ поправкой ранга два:
\[
    \vec{s}_k = \vec{\theta}^{\,k+1} - \vec{\theta}^{\,k} = \alpha_k \vec{p}_k,
    \qquad
    \vec{y}_k = \nabla f(\vec{\theta}^{\,k+1}) - \nabla f(\vec{\theta}^{\,k}),
    \qquad
    \rho_k = \frac{1}{\vec{y}_k^{\,T} \vec{s}_k},
\]
\[
    B_{k+1} = (I - \rho_k\, \vec{s}_k \vec{y}_k^{\,T})\, B_k\, (I - \rho_k\, \vec{y}_k \vec{s}_k^{\,T}) + \rho_k\, \vec{s}_k \vec{s}_k^{\,T},
\]
направление спуска $\vec{p}_k = -B_k \nabla f(\vec{\theta}^{\,k})$, шаг $\alpha_k$ подбирается \texttt{LineSearch}; при нарушении положительной определённости матрица сбрасывается в единичную.

\subsectiontext{4. Методы нулевого порядка}

\textbf{Метод Нелдера–Мида} (\texttt{NelderMead}). В $n$-мерном пространстве строится симплекс из $n+1$ вершины $\{\vec{x}_1, \dots, \vec{x}_{n+1}\}$; вершины сортируются по значениям $f$, отражение выполняется наихудшей вершины через центроид $\vec{x}_o$ лучших $n$ вершин, $\vec{x}_r = \vec{x}_o + \alpha (\vec{x}_o - \vec{x}_{n+1})$, $\alpha = 1$. Далее применяются стандартные операции с коэффициентами $\gamma = 2$ (растяжение $\vec{x}_e = \vec{x}_o + \gamma(\vec{x}_r - \vec{x}_o)$), $\rho = 0{,}5$ (внешнее и внутреннее сжатие) и $\sigma = 0{,}5$ (редукция — стягивание всех вершин к лучшей). Критерий останова — разброс расстояний вершин до центроида:
\[
    \max_i \|\vec{x}_i - \vec{x}_c\| < \varepsilon, \qquad \vec{x}_c = \frac{1}{n+1}\sum_i \vec{x}_i.
\]

\textbf{Метод роя частиц (PSO)} (\texttt{PSO}). Скорость и позиция $j$-й частицы обновляются по схеме
\[
    \vec{v}_j \leftarrow \omega_k\, \vec{v}_j + c_{1k}\, r_1 \odot (\vec{p}_j^{\,best} - \vec{x}_j) + c_{2k}\, r_2 \odot (\vec{g}^{\,best} - \vec{x}_j),
    \qquad
    \vec{x}_j \leftarrow \vec{x}_j + \vec{v}_j,
\]
где $r_1, r_2 \sim U(0,1)$. Реализована модифицированная схема с адаптивными коэффициентами: инерция линейно убывает $\omega_k = \omega_{max} - (\omega_{max} - \omega_{min})\, k / k_{max}$, $\omega_{max} = 0{,}9$, $\omega_{min} = 0{,}4$; когнитивный и социальный веса каждой частицы зависят от отставания её локального рекорда $F^{LB}$ от глобального $F^{GB}$:
\[
    c_{1k} = c_{min} + (c_{max} - c_{min})\, \gamma_k,
    \qquad
    c_{2k} = c_{max} - (c_{max} - c_{min})\, \gamma_k,
    \qquad
    \gamma_k = \frac{F^{LB} - F^{GB}}{\max_j (F^{LB}_j - F^{GB})},
\]
то есть отставая от лидера роя, частица сильнее опирается на собственный рекорд, а приближаясь к нему — на социальную тягу ($c_{min} = 0{,}5$, $c_{max} = 2{,}5$).

\subsectiontext{5. Нелинейный метод наименьших квадратов}

Задача ставится как минимизация суммы квадратов невязок $S(\vec{\theta}) = \|\vec{r}(\vec{\theta})\|^2$ с градиентом $\nabla S = J^T \vec{r}$.

\textbf{Метод Гаусса–Ньютона} (\texttt{GaussNewton}) аппроксимирует гессиан $H \approx J^T J$ и решает на каждой итерации систему
\[
    (J^T J)\, \vec{h} = -J^T \vec{r},
\]
при необходимости дробя шаг вдвое до убывания невязки.

\textbf{Метод Левенберга–Марквардта} (\texttt{LevenbergMarquardt}) добавляет адаптивное демпфирование:
\[
    \bigl(J^T J + \lambda\, \operatorname{diag}(J^T J)\bigr)\, \vec{h}_{lm} = -J^T \vec{r}.
\]
Качество шага оценивается метрикой выигрыша
\[
    \rho(\vec{h}_{lm}) = \frac{S(\vec{\theta}) - S(\vec{\theta} + \vec{h}_{lm})}{\vec{h}_{lm}^{\,T} J^T J\, \vec{h}_{lm} + 2\lambda\, \vec{h}_{lm}^{\,T} \operatorname{diag}(J^T J)\, \vec{h}_{lm}};
\]
при $\rho > 0$ шаг принимается и $\lambda$ уменьшается в $\nu$ раз, иначе $\lambda$ увеличивается в $\nu$ раз и шаг пересчитывается ($\nu = 10$).

\subsectiontext{6. Прикладная задача с ограничениями: распределение времени на сессию}

Максимизируется суммарная вероятность сдачи $N = 3$ экзаменов при бюджете времени $M = 8$ ч:
\[
    \max_{\vec{t}}\ \sum_{i=1}^{3} P_i(t_i)
    \quad \text{при} \quad
    \sum_{i=1}^{3} t_i \le M,\ \ t_i > 0,
\]
что сводится к минимизации $f(\vec{t}) = -\sum_i P_i(t_i)$. Первые два экзамена моделируются логистическими кривыми $P_i(t) = 1/(1 + e^{-k_i(t - t_{0i})})$ с параметрами $(k_1, t_{01}) = (1{,}5;\, 3{,}0)$, $(k_2, t_{02}) = (2{,}0;\, 4{,}5)$, третий — функцией типа Гаусса $P_3(t) = \exp\bigl(-(t-2)^2 / 1{,}5\bigr)$.

Ограничения учитываются двумя способами.

\textbf{Проекция на допустимое множество} (класс \texttt{ProjectedGradient}): шаг градиентного спуска с \texttt{LineSearch} чередуется с евклидовой проекцией на симплекс $\{\vec{t} \ge 0: \sum_i t_i = M\}$ (функция \texttt{project\_simplex}); при необходимости шаг дробится до убывания $f$. Остановка — по невязке условий Куна–Таккера $\|\vec{x} - \Pi(\vec{x} - \nabla f)\| < \varepsilon$.

\textbf{Метод штрафных функций} (класс \texttt{QuadraticPenalty}): минимизация $F_\mu(\vec{t}) = f(\vec{t}) + \mu\, h(\vec{t})$ методом BFGS с квадратичным штрафом
\[
    h(\vec{t}) = \sum_{i=1}^{N} \max^2(0, -t_i) + \max^2\Bigl(0, \sum_{i=1}^{N} t_i - M\Bigr),
\]
параметр $\mu$ наращивается по схеме $10 \to 100 \to 1000$.

Для анализа устойчивости выполняется серия из 100 запусков обоих подходов со случайных начальных приближений $t_i \sim U(0{,}1;\, M)$; сошедшиеся решения группируются в кластеры округлением координат до двух знаков, для каждого кластера вычисляются средняя точка, число траекторий и значение $f$.

\newpage

% --- Практическая часть ---
\label{sec:practice}
\mainsection{Практическая часть}

Все расчёты выполняются одной командой \texttt{python main.py}, которая последовательно запускает пять блоков и сохраняет графики в файлы \texttt{png}. Для каждого метода фиксируется траектория приближений и печатается сводка: число итераций, точка окончания и значение функции.

\subsectiontext{1. Градиентные методы первого порядка}

Методы тестируются на двумерных функциях Розенброка (старт $(-1{,}5;\, 2{,}0)$) и Химмельблау (старт $(-4;\, 4)$); каждый метод прогоняется с аналитическим и численным градиентом. Результаты сведены в таблицы~\ref{tab:block1} и~\ref{tab:block1h}, траектории показаны на рис.~\ref{fig:rosenbrock} и~\ref{fig:himmelblau}.

\begin{table}[H]
    \centering
    \small
    \caption{Сходимость методов первого порядка на функции Розенброка}
    \label{tab:block1}
    \begin{tabular}{llrrl}
        \toprule
        Метод & Градиент & Итераций & $f(\vec{\theta}^{\,*})$ & $\vec{\theta}^{\,*}$ \\
        \midrule
        Наискорейший спуск & аналит. & 18761 & $6.4\cdot10^{-11}$ & $(1.00001;\ 1.00002)$ \\
                           & численный & 16369 & $1.2\cdot10^{-10}$ & $(1.00001;\ 1.00002)$ \\
        Градиентный спуск (фикс. шаг) & аналит. & 10027 & $1.2\cdot10^{-10}$ & $(0.99999;\ 0.99998)$ \\
                           & численный & 10002 & $1.3\cdot10^{-10}$ & $(0.99999;\ 0.99998)$ \\
        Покомпонентный ($x \to y$) & аналит. & 9804 & $2.5\cdot10^{-11}$ & $(1.00000;\ 1.00001)$ \\
                           & численный & 9806 & $2.5\cdot10^{-11}$ & $(1.00000;\ 1.00001)$ \\
        Покомпонентный (перемеш.) & аналит. & 12318 & $9.8\cdot10^{-11}$ & $(1.00001;\ 1.00002)$ \\
                           & численный & 12282 & $1.0\cdot10^{-10}$ & $(1.00001;\ 1.00002)$ \\
        Покомпонентный (фикс. шаг) & аналит. & 18486 & $1.2\cdot10^{-10}$ & $(0.99999;\ 0.99998)$ \\
                           & численный & 18486 & $1.2\cdot10^{-10}$ & $(0.99999;\ 0.99998)$ \\
        Момент Нестерова & аналит. & 1176 & $1.2\cdot10^{-10}$ & $(0.99999;\ 0.99998)$ \\
                           & численный & 1176 & $1.2\cdot10^{-10}$ & $(0.99999;\ 0.99998)$ \\
        Adam & аналит. & 1840 & $9.5\cdot10^{-11}$ & $(1.00000;\ 0.99998)$ \\
                           & численный & 1840 & $9.6\cdot10^{-11}$ & $(1.00000;\ 0.99998)$ \\
        Adagrad & аналит. & 19863 & $1.2\cdot10^{-10}$ & $(0.99999;\ 0.99998)$ \\
                           & численный & 19863 & $1.2\cdot10^{-10}$ & $(0.99999;\ 0.99998)$ \\
        RMSProp & аналит. & 20000 & $2.2\cdot10^{-2}$ & $(0.990;\ 0.966)$ \\
                           & численный & 20000 & $2.2\cdot10^{-2}$ & $(0.980;\ 0.976)$ \\
        \bottomrule
    \end{tabular}
\end{table}

\begin{table}[H]
    \centering
    \small
    \caption{Сходимость методов первого порядка на функции Химмельблау}
    \label{tab:block1h}
    \begin{tabular}{llrrl}
        \toprule
        Метод & Градиент & Итераций & $f(\vec{\theta}^{\,*})$ & $\vec{\theta}^{\,*}$ \\
        \midrule
        Наискорейший спуск & аналит. & 27 & $4.1\cdot10^{-13}$ & $(3.584;\ -1.848)$ \\
                           & численный & 27 & $4.1\cdot10^{-13}$ & $(3.584;\ -1.848)$ \\
        Градиентный спуск (фикс. шаг) & аналит. & 111 & $6.3\cdot10^{-13}$ & $(-2.805;\ 3.131)$ \\
                           & численный & 111 & $6.3\cdot10^{-13}$ & $(-2.805;\ 3.131)$ \\
        Покомпонентный ($x \to y$) & аналит. & 16 & $1.5\cdot10^{-14}$ & $(-3.779;\ -3.283)$ \\
                           & численный & 16 & $1.5\cdot10^{-14}$ & $(-3.779;\ -3.283)$ \\
        Покомпонентный (перемеш.) & аналит. & 20 & $1.9\cdot10^{-13}$ & $(-3.779;\ -3.283)$ \\
                           & численный & 18 & $1.9\cdot10^{-13}$ & $(-3.779;\ -3.283)$ \\
        Покомпонентный (фикс. шаг) & аналит. & 222 & $6.3\cdot10^{-13}$ & $(-2.805;\ 3.131)$ \\
                           & численный & 222 & $6.3\cdot10^{-13}$ & $(-2.805;\ 3.131)$ \\
        Момент Нестерова & аналит. & 241 & $2.5\cdot10^{-13}$ & $(-2.805;\ 3.131)$ \\
                           & численный & 241 & $2.5\cdot10^{-13}$ & $(-2.805;\ 3.131)$ \\
        Adam & аналит. & 290 & $2.7\cdot10^{-13}$ & $(-2.805;\ 3.131)$ \\
                           & численный & 290 & $2.7\cdot10^{-13}$ & $(-2.805;\ 3.131)$ \\
        Adagrad & аналит. & 23 & $2.4\cdot10^{-13}$ & $(-2.805;\ 3.131)$ \\
                           & численный & 23 & $2.4\cdot10^{-13}$ & $(-2.805;\ 3.131)$ \\
        RMSProp & аналит. & 133 & $6.2\cdot10^{-13}$ & $(-2.805;\ 3.131)$ \\
                           & численный & 133 & $6.2\cdot10^{-13}$ & $(-2.805;\ 3.131)$ \\
        \bottomrule
    \end{tabular}
\end{table}

На функции Розенброка (рис.~\ref{fig:rosenbrock}) все методы, кроме RMSProp, достигают глобального минимума $(1; 1)$ с точностью не хуже $10^{-10}$. Наискорейший спуск из-за точного одномерного поиска делает характерные зигзагообразные микродвижения по дну изогнутой долины и требует больше всего итераций (более 18 тысяч). Его вариант с фиксированным шагом $\text{lr} = 2\cdot10^{-3}$ сходится почти вдвое быстрее (10 тысяч итераций) и осцилляции у дна долины затухают, что видно на траектории: шаг, безопасный для всей траектории ($L \approx 10^3$ в долине), заметно консервативнее одномерного поиска, зато не требует дополнительных вычислений $f$. Покомпонентный спуск движется ломаной с прямыми участками вдоль осей (видны оба варианта обхода); его вариант с фиксированным шагом ведёт себя аналогично градиентному, но в полтора раза медленнее из-за покомпонентного характера движения. Момент Нестерова и Adam, использующие накопленную «инерцию» и адаптивный шаг, проходят долину уверенно и сходятся на порядок быстрее. Adagrad также достигает минимума, однако монотонно растущая сумма $G_t$ постепенно глушит шаг, и метод по числу итераций догоняет наискорейший спуск. RMSProp со скользящим средним удерживает рабочий размер шага и быстро спускается в долину, но в её дне, где $\vec{v}_t \approx (\nabla f_t)^2$, итерации вырождаются в блуждание с амплитудой $\sim \text{lr}$ вокруг минимума: за 20 тысяч итераций точность по норме градиента $10^{-5}$ не достигается ($f \approx 2.2\cdot10^{-2}$). Это ожидаемое поведение постоянношаговой версии метода на изогнутой долине. Численный градиент даёт траектории, практически неотличимые от аналитического, что подтверждает корректность конечно-разностной схемы.

На функции Химмельблау (рис.~\ref{fig:himmelblau}) из точки $(-4; 4)$ разные методы сходятся к разным минимумам, что наглядно демонстрирует многомодальность функции: наискорейший спуск за 27 итераций приходит в $(3.584;\, -1.848)$, оба покомпонентных метода со случайным и последовательным обходом движутся ломаными вдоль осей и завершаются в $(-3.779;\, -3.283)$, метод Нестерова, Adam, Adagrad, RMSProp и оба спуска с фиксированным шагом приходят в минимум $(-2.805;\, 3.131)$, попадая в его бассейн притяжения. Adagrad оказывается быстрее всех (23 итерации) — на хорошо обусловленных участках его затухающий шаг после первых резких шагов точно доводит траекторию до минимума; спуски с фиксированным шагом сходятся устойчиво, но уступают в числе итераций. Заметим, что методы с фиксированным шагом (Нестеров, Adam, Adagrad, RMSProp и оба спуска с фиксированным множителем) ведут себя стабильно на многомодальной функции, тогда как на унимодальной Розенброке чувствительность к выбору $\text{lr}$ проявляется сильнее.

\begin{figure}[H]
    \centering
    \img[width=0.95\textwidth]{rosenbrock}
    \caption{Траектории методов первого порядка на функции Розенброка}
    \label{fig:rosenbrock}
\end{figure}

\begin{figure}[H]
    \centering
    \img[width=0.95\textwidth]{himmelblau}
    \caption{Траектории методов первого порядка на функции Химмельблау}
    \label{fig:himmelblau}
\end{figure}

\subsectiontext{2. Методы второго порядка}

Метод Ньютона прогоняется с аналитическим и численным гессианом, BFGS строит аппроксимацию обратного гессиана самостоятельно. Результаты — в таблице~\ref{tab:block2}, траектории — на рис.~\ref{fig:second}.

\begin{table}[H]
    \centering
    \small
    \caption{Сходимость методов второго порядка}
    \label{tab:block2}
    \begin{tabular}{llrrl}
        \toprule
        Метод & Функция & Итераций & $f(\vec{\theta}^{\,*})$ & $\vec{\theta}^{\,*}$ \\
        \midrule
        Ньютон (аналит. $H$) & Розенброк & 23 & $2.5\cdot10^{-28}$ & $(1;\ 1)$ \\
        Ньютон (числ. $H$) & Розенброк & 23 & $3.0\cdot10^{-27}$ & $(1;\ 1)$ \\
        BFGS & Розенброк & 17 & $4.3\cdot10^{-18}$ & $(1;\ 1)$ \\
        \midrule
        Ньютон (аналит. $H$) & Химмельблау & 5 & $1.3\cdot10^{-22}$ & $(-2.805;\ 3.131)$ \\
        Ньютон (числ. $H$) & Химмельблау & 5 & $1.3\cdot10^{-22}$ & $(-2.805;\ 3.131)$ \\
        BFGS & Химмельблау & 8 & $2.2\cdot10^{-18}$ & $(3.584;\ -1.848)$ \\
        \bottomrule
    \end{tabular}
\end{table}

Методы второго порядка заметно экономнее первого: Ньютон сходится за считанные итерации, его траектория практически не чувствует изогнутости долины Розенброка и скользит вдоль неё к точке $(1; 1)$, дав точность до $10^{-28}$. Численный гессиан (якобиан численного градиента) даёт ту же траекторию и то же число итераций. BFGS на первых итерациях ведёт себя как наискорейший спуск ($B_0 = I$), затем уточнение $B_k$ позволяет ему развернуться вдоль долины и сойтись за 17 итераций — быстрее всех методов первого порядка. На Химмельблау методы из одной стартовой точки $(-4; 4)$ попадают в разные бассейны притяжения: Ньютон — в минимум $(-2.805;\, 3.131)$ всего за 5 итераций, BFGS — в $(3.584;\, -1.848)$ за 8 итераций, что ещё раз иллюстрирует многомодальность функции.

\begin{figure}[H]
    \centering
    \img[width=0.95\textwidth]{second_order}
    \caption{Траектории методов второго порядка: Ньютон (аналитический и численный гессиан) и BFGS}
    \label{fig:second}
\end{figure}

\subsectiontext{3. Методы нулевого порядка}

\textbf{Нелдер–Мид.} На двумерной функции Розенброка симплекс из трёх вершин стартует из точки $(-1{,}5; 2{,}0)$; на рис.~\ref{fig:nm2d} белыми линиями показано положение симплекса в характерные моменты, красной — траектория его центроида. Метод сперва стягивает симплекс и перетаскивает его по дну долины, затем ведёт к минимуму $(1; 1)$ за 128 итераций с $f \approx 5.2\cdot10^{-18}$. На шестимерной функции (старт $(-1, \dots, -1)^T$, рис.~\ref{fig:nm6d}) наблюдается классическая для Розенброка картина: длительные плато медленного спуска (между 50-й и 450-й итерациями $f$ убывает всего с $10^{0}$ до $10^{-1}$), после чего сходимость резко ускоряется и за $\approx 200$ последних итераций невязка падает ниже $10^{-13}$; всего потребовалось 673 итерации, $f^* \approx 1.2\cdot10^{-14}$.

\begin{figure}[H]
    \centering
    \img[width=0.8\textwidth]{nelder_mead_2d}
    \caption{Нелдер–Мид: траектория центроида и положение симплекса (2D Розенброк)}
    \label{fig:nm2d}
\end{figure}

\begin{figure}[H]
    \centering
    \img[width=0.8\textwidth]{nelder_mead_6d}
    \caption{Нелдер–Мид: сходимость $f(\vec{\theta}^{\,*_k})$ на шестимерной функции Розенброка}
    \label{fig:nm6d}
\end{figure}

\textbf{Рой частиц (PSO).} Рой из 30 частиц минимизирует двумерную функцию Розенброка в области $[-2; 2] \times [-1; 3]$ в течение 150 итераций; на рис.~\ref{fig:pso} показаны позиции частиц и векторы их скоростей на итерациях 0, 50, 100 и 150. В начале частицы равномерно рассеяны по области (итерация 0); к 50-й итерации рой заметно концентрируется в долине, белая звезда — текущий глобальный лидер; к 100-й итерации большинство частиц локализовано в окрестности минимума, скорости ещё значимы; к 150-й итерации рой схлопывается в точке $(1; 1)$. Адаптивные коэффициенты обеспечивают баланс глобального поиска на старте и локальной сходимости в конце: вместе с убывающей инерцией они позволяют рою пройти узкую долину и найти глобальный минимум с $f \approx 2.2\cdot10^{-10}$ в точке $(1.00000;\, 1.00000)$ без единого вычисления производных. Результаты блока сведены в таблицу~\ref{tab:block3}.

\begin{figure}[H]
    \centering
    \img[width=0.85\textwidth]{pso}
    \caption{PSO: распределение частиц и их скоростей на изолиниях Розенброка}
    \label{fig:pso}
\end{figure}

\begin{table}[H]
    \centering
    \small
    \caption{Результаты методов нулевого порядка}
    \label{tab:block3}
    \begin{tabular}{llrrl}
        \toprule
        Метод & Задача & Итераций & $f(\vec{\theta}^{\,*})$ & $\vec{\theta}^{\,*}$ \\
        \midrule
        Нелдер–Мид & Розенброк 2D & 128 & $5.2\cdot10^{-18}$ & $(1;\ 1)$ \\
        Нелдер–Мид & Розенброк 6D & 673 & $1.2\cdot10^{-14}$ & $f(\theta^{*_k})$, рис.~\ref{fig:nm6d} \\
        PSO (30 частиц) & Розенброк 2D & 150 & $2.2\cdot10^{-10}$ & $(1.00000;\ 1.00000)$ \\
        \bottomrule
    \end{tabular}
\end{table}

\subsectiontext{4. Нелинейный метод наименьших квадратов}

Функция Розенброка записывается как задача МНК с невязками $r(\vec{\theta}) = (1 - \theta_1;\ 10(\theta_2 - \theta_1^2))^T$; оба метода стартуют из $(-1{,}5; 2{,}0)$. Результаты — в таблице~\ref{tab:block4}, траектории и графики — на рис.~\ref{fig:ls}. Оба метода проходят долину и сходятся к $(1; 1)$; Левенберг–Марквардт делает это плавнее за счёт демпфирования, Гаусс–Ньютон использует дробление шага.

В задаче аппроксимации генерируются $n = 40$ зашумлённых точек $(x_i, y_i)$ истинной моделью
\[
    g_{\vec{\theta}}(x) = a\, e^{-((x-b)/c)^2} + d\, x
\]
с параметрами $\vec{\theta}_{true} = (2{,}0;\, 1{,}5;\, 0{,}8;\, 0{,}3)^T$ и шумом $\xi \sim N(0; 0{,}05)$; невязки $r_i(\vec{\theta}) = g_{\vec{\theta}}(x_i) - y_i$, якобиан вычисляется конечными разностями. Нижний ряд рис.~\ref{fig:ls} показывает эволюцию кривой $g_{\vec{\theta}}(x)$ по итерациям Левенберга–Марквардта (начальная прямая сгибается в гауссову кривую и подгоняется под данные уже к 5-й итерации) и кривые сходимости невязки: оба метода снижают $S(\vec{\theta}_k)$ с $\approx 40$ до $5.1\cdot10^{-2}$ — уровня, определяемого шумом данных. Левенберг–Марквардт сходится за 11 итераций, а Гаусс–Ньютон, чей полный шаг на старте далёк от минимума, работает до ограничения в 300 итераций, но приходит в ту же точку $\vec{\theta}^{\,*} \approx (1.966;\, 1.507;\, 0.811;\, 0.300)$, близкую к истинной. Результаты — в таблице~\ref{tab:block4}.

\begin{figure}[H]
    \centering
    \img[width=0.95\textwidth]{least_squares}
    \caption{МНК: траектории на Розенброке, эволюция кривой аппроксимации и сходимость невязки}
    \label{fig:ls}
\end{figure}

\begin{table}[H]
    \centering
    \small
    \caption{Результаты методов нелинейных наименьших квадратов}
    \label{tab:block4}
    \begin{tabular}{llrl}
        \toprule
        Метод & Задача & Итераций & $f(\vec{\theta}^{\,*})$ \\
        \midrule
        Гаусс–Ньютон & Розенброк 2D & 15 & $5.0\cdot10^{-30}$ \\
        Левенберг–Марквардт & Розенброк 2D & 23 & $9.5\cdot10^{-24}$ \\
        Гаусс–Ньютон & аппроксимация данных & 300 (предел) & $5.1\cdot10^{-2}$ \\
        Левенберг–Марквардт & аппроксимация данных & 11 & $5.1\cdot10^{-2}$ \\
        \bottomrule
    \end{tabular}
\end{table}

\subsectiontext{5. Распределение времени на сессию}

Выполнено по $n = 100$ запусков обоих подходов со случайных начальных приближений $t_i \sim U(0{,}1; 8)$. Итоговые кластеры сошедшихся решений приведены в таблице~\ref{tab:session}.

\begin{table}[H]
    \centering
    \small
    \caption{Мультизапуск (100 стартов): кластеры решений задачи распределения времени}
    \label{tab:session}
    \begin{tabular}{llrl}
        \toprule
        Подход & $(t_1;\ t_2;\ t_3)$, ч & Траекторий & $F(\vec{t})$ \\
        \midrule
        Проекция на симплекс & $(6.01;\ 0.00;\ 1.99)$ & 69 & $-1.989$ \\
        Проекция на симплекс & $(0.00;\ 6.06;\ 1.94)$ & 24 & $-1.966$ \\
        Проекция на симплекс & $(2.95;\ 5.05;\ 0.00)$ & 7 & $-1.301$ \\
        \midrule
        Квадратичный штраф (BFGS) & $(6.01;\ 0.00;\ 1.99)$ & 75 & $-1.989$ \\
        Квадратичный штраф (BFGS) & $(0.00;\ 6.06;\ 1.94)$ & 25 & $-1.966$ \\
        \bottomrule
    \end{tabular}
\end{table}

Мультизапуск показывает, что задача многомодальна: оба метода находят два близких по качеству решения, а доминирующим (69–75 траекторий из 100) является точка $(6.01;\, 0;\, 1.99)$ с $F \approx -1.989$. Физический смысл результата таков: почти всё время отдаётся первому экзамену, чья логистическая кривая при $t_1 \approx 6$ практически насыщается ($P_1 \approx 0.99$), и третьему — ровно столько, чтобы встать в пик его гауссианы ($t_3 \approx 2$, $P_3 \approx 1$). Второму экзамену время не выделяется: его сигмоида центрирована при $t_{02} = 4{,}5$, поэтому при $t_2 = 0$ вероятность почти нулевая ($P_2 \approx 10^{-4}$), а подъём её до заметных значений потребовал бы слишком большой доли бюджета. Проверка предельных выигрышей подтверждает оптимум: в точке решения производные $\partial P_1 / \partial t_1 \approx 0.016$ и $\partial P_3 / \partial t_3 \approx 0.013$ выровнены, тогда как $\partial P_2 / \partial t_2 \approx 2.5\cdot10^{-4}$ — почти в 60 раз меньше.

Второй кластер $(0;\, 6.06;\, 1.94)$ соответствует альтернативной стратегии «спасти второй экзамен»: часть вероятности первого жертвуется ради частичного насыщения второго, итоговая суммарная вероятность чуть ниже ($-1.966$). Метод проекции из 100 стартов семь раз попал ещё и в худший локальный минимум $(2.95;\, 5.05;\, 0)$ с $F = -1.301$ (делёж времени между первым и вторым экзаменами без учёта третьего); штрафной метод с теми же стартами в него не попал. Доминирующие кластеры обоих подходов совпадают почти полностью (69 против 75 траекторий), а значения $F$ согласованы с точностью до тысячных, что подтверждает эквивалентность проекции и штрафного метода, несмотря на принципиально разные механизмы учёта ограничений.

\newpage

% --- Вывод ---
\label{sec:conclusion}
\mainsection{Вывод}

В лабораторной работе реализован единый программный комплекс численной оптимизации на языке Python с модульной архитектурой: базовый класс \texttt{TargetFunction} с аналитическим и конечно-разностным дифференцированием, базовый класс \texttt{Optimizer} с единообразным накоплением истории и метрик сходимости, класс \texttt{LineSearch} (золотое сечение) для одномерного поиска шага, а также самостоятельные модули для методов нулевого порядка, нелинейных наименьших квадратов и задачи с ограничениями.

Реализовано и протестировано пятнадцать методов: градиентные первого порядка (наискорейший спуск, его вариант с фиксированным шагом, покомпонентный спуск в двух вариантах обхода и с фиксированным шагом, момент Нестерова, Adam, Adagrad, RMSProp), второго порядка (Ньютон, BFGS), нулевого порядка (Нелдер–Мид, PSO с адаптивными коэффициентами), методы Гаусса–Ньютона и Левенберга–Марквардта, а также два способа учёта ограничений — проекция на симплекс и квадратичная штрафная функция.

Вычислительные эксперименты на функциях Розенброка и Химмельблау показали: сравнение аналитического и численного дифференцирования подтвердило корректность обеих схем (траектории практически совпадают); спуск с фиксированным шагом уступает наискорейшему по надёжности и сильно зависит от обусловленности задачи, а адаптивные методы (Adam, Adagrad, RMSProp) проходят изогнутую долину Розенброка без одномерного поиска — каждый со своими особенностями: у Adagrad затухает шаг, у RMSProp с постоянным $\text{lr}$ устанавливается блуждание окрестности минимума; методы второго порядка сходятся за считанные итерации, но требуют гессиана или его аппроксимации; безпроизводные методы успешно работают в многомерном случае (Нелдер–Мид на 6-мерной Розенброке), а PSO с адаптивной инерцией и весами находит глобальный минимум без вычисления производных. Методы МНК уверенно решают как минимизацию Розенброка, так и аппроксимацию зашумлённых данных четырёхпараметрической моделью, снижая невязку до уровня шума.

В прикладной задаче распределения времени на сессию оба подхода к ограничениям (проекция на симплекс и квадратичный штраф) при мультизапуске из 100 случайных точек сходятся к согласованному решению $(6.01;\, 0;\, 1.99)$ часов с суммарной вероятностью сдачи $\approx 1.99$, рационально распределяющему бюджет $M = 8$ часов между экзаменами и подтверждающему устойчивость найденного оптимума.

\end{document}
